% !TeX root = ../main.tex

\section{Threats}
% We summarize three threats to our evaluation. First, the number of malicious data exfiltration packages is limited. However, we mitigated this by a random sampling to select representative subset and cover diverse behaviors and patterns. Second, the manual analysis may introduce bias, as subjective judgment is involved in labeling and interpreting attack behaviors and patterns. We followed a consistent methodology and cross-validated results confirmed by the three of the authors to reduce potential bias. Third, the evaluation relies on specific benchmark settings and tool. However, these settings were chosen to reflect realistic scenarios as closely as possible.
\todo{We summarize five threats to our evaluation. First, the number of malicious data exfiltration packages is limited. However, we mitigated this by a random sampling to select representative subset and cover diverse behaviors and patterns. Second, the manual analysis may introduce bias, as subjective judgment is involved in labeling and interpreting attack behaviors and patterns. We followed a consistent methodology and cross-validated results confirmed by the three of the authors to reduce potential bias. Third, the evaluation treats detection tools as black-box systems using their default configurations, which may not fully reflect the maximum capability of each tool. However, this design reflects realistic usage scenarios where users typically rely on default settings. Fourth, some of the evaluated tools relied on sandbox-based dynamic analysis, where certain malicious behaviors may depend on runtime conditions, which may lead to underestimation of performance by dynamic analysis tools. To mitigate this limitation, we also included static-analysis–based detectors, allowing the evaluation to cover complementary detection paradigms. Finally, the evaluation is conducted under specific benchmark settings, including predefined dilution ratios and detection tools. However, these settings were designed to approximate realistic package ecosystem conditions and to examine tool behavior under different class imbalance scenarios.}